% WELCOME -- the notes' sections 0.1 and 0.4, sec.theCourse and sec.learningFinance.
% The first frames of the course. Budget 4 min.
% sec.learningFinance has no frame of its own: what Chapter 1 builds is the second and third
% rows of the table below, where it costs no words. The right-hand column is the part of
% the package that the laboratory half of the course works on.
%
% THE CAPTION CONTRACT. A caption is a gloss (`term: noun phrase'), a fragment of subject
% and predicate, or a condition written as mathematics -- never a sentence, and never a
% repetition of the frame title. At most 12 words each, three to a frame, 25 to a frame in
% total. What is said over the object goes in \note{}. Check with
%   python3 .claude/skills/writing-in-tex/scripts/frame_density.py <file>

\subsection{The course}

\begin{frame}
  \frametitle{Market microstructure: lectures, then the lab}
\begin{center}
\small
\begin{tabular}{@{}lll@{}}
\toprule
 & \emph{in lectures} & \emph{in the lab} \\
 & Chapter 1 & its implementation \\
\midrule
the subject & orders, matching, the book & the modules that hold them \\
we build & the fold, and faster folds & five books, one matching \\
on the way & Hawkes processes & two simulation schemes \\
the end point & price formation & sessions from recorded data \\
\bottomrule
\end{tabular}
\end{center}

In Python. Every computable statement, computed.

Your questions, by email: what the lab concentrates on.

Chapter 2, options and volatility: optional, not examined.

\url{https://github.com/claudio-ICL/unito26}
\note{
The course is about market microstructure, Chapter 1 of the notes: what an order is, how an
exchange matches orders, what the limit order book is at any instant, and how the flow of
orders forms a price. We fold the stream of orders into an aggregated book, examine several
implementations of that fold and compare their performance, then model the arrival of orders
as a Hawkes process and ask what its parametrisation says about a trading epoch. The order
flow imbalance and price formation are the end point.
The first half of the course presents the chapter in lectures. The second half is laboratory
work on its implementation: we read the code, run it and change it together. Write to me with
questions and remarks on what you find challenging in the implementation, and the laboratory
sessions concentrate on those parts.
Chapter 2 of the notes is optional reading on option pricing and volatility trading. It is
neither taught in class nor examined, and it is added to the notes towards the end of the
course.
You are expected to read and understand the implementations, not only to run them.
}
\end{frame}

\begin{frame}
  \frametitle{Writing code that works no longer gates the analysis}
\begin{center}
\begin{tikzpicture}[font=\scriptsize, node distance=6mm,
    box/.style={draw=gray!60, rounded corners=2pt, align=center, inner sep=4pt,
                text width=2.0cm, fill=gray!8},
    tool/.style={box, fill=blue!12, draw=blue!50},
    keep/.style={draw=orange!70, rounded corners=2pt, align=left, inner sep=5pt,
                 text width=3.7cm, fill=orange!15}]
  \node[box] (code) {code that works};
  \node[box, right=10mm of code] (analysis) {the analysis};
  \node[tool, below=of code] (llm) {coding assistants};
  \node[keep, right=6mm of analysis] (judgment)
       {\emph{judgment}\\ correct, or merely running\\ two designs, one choice\\
        the cost a year later};
  \draw[-{Stealth}, gray!60] (code) -- node[above, font=\tiny, gray] {the gate} (analysis);
  \draw[-{Stealth}, blue!50] (llm) -- (code);
  \draw[-{Stealth}, orange!70] (judgment) -- (analysis);
\end{tikzpicture}
\end{center}

\note{
Until recently, writing code that worked was the skill that gated everything else in
quantitative analysis. With coding assistants that is no longer true.
What an assistant does not supply is the judgment to tell a correct implementation from one
that merely runs, to choose between two designs, and to see what an implementation choice
will cost a year later.
That judgment is what this course teaches, and it is what the examination tests.
}
\end{frame}
